% Compatibilidad con instalaciones MiKTeX cuyo formato base es anterior a
% la versión instalada de Hyperref.
\providecommand{\IfFormatAtLeastTF}[3]{#3}
\providecommand{\IfFormatAtLeastT}[2]{}
\providecommand{\IfFormatAtLeastF}[2]{#2}
\providecommand{\IfDocumentMetadataTF}[2]{#2}
\providecommand{\IfDocumentMetadataT}[1]{}
\providecommand{\IfDocumentMetadataF}[1]{#1}

\documentclass[11pt]{article}

\usepackage[T1]{fontenc}
\usepackage[utf8]{inputenc}
\usepackage[spanish,es-nodecimaldot]{babel}
\usepackage{lmodern}
\usepackage[margin=1in]{geometry}
\usepackage{amsmath,amssymb}

\setlength{\parindent}{0pt}
\setlength{\parskip}{0.45em}
\renewcommand{\labelenumi}{\textnormal{(\alph{enumi})}}
\pagestyle{plain}

\let\enumerateoriginal\enumerate
\let\endenumerateoriginal\endenumerate
\renewenvironment{enumerate}{%
  \enumerateoriginal
  \setlength{\itemsep}{0.5em}%
  \setlength{\parsep}{0pt}%
  \setlength{\topsep}{0.4em}%
}{\endenumerateoriginal}

\newcounter{problema}
\newenvironment{problema}[1]{%
  \par
  \refstepcounter{problema}%
  \addvspace{0.9em}%
  \noindent\textbf{Ejercicio \theproblema. #1}
  \par\nobreak\smallskip
}{\par}

\begin{document}

\begin{center}
  {\large\textbf{Centro de Investigación y Docencia Económicas}}\\[0.25em]
  {\large\textbf{Macroeconomía II}}\\[0.2em]
  {\large\textbf{Problem Set 3: Consumo intertemporal y mercados de crédito}}\\[0.35em]
  LECO -- Otoño 2026
\end{center}

\vspace{0.6em}
\hrule
\vspace{0.6em}



\begin{problema}{Cambios transitorios y permanentes en los impuestos}
Considere un hogar con preferencias estrictamente crecientes y estrictamente convexas sobre $(c,c')$, que recibe ingresos $(y,y')$ y paga impuestos de suma fija $(\tau,\tau')$. Puede ahorrar o endeudarse libremente a la tasa de interés real $r$. Suponga que ambos consumos son bienes normales.

Para cada uno de los siguientes casos, utilice un diagrama en el plano $(c,c')$ para analizar el efecto sobre el consumo presente, el consumo futuro y el ahorro actual.

\begin{enumerate}
  \item El gobierno anuncia un aumento de los impuestos en $\Delta\tau>0$ únicamente en el periodo actual.

  \item El gobierno anuncia un aumento de impuestos permanente: tanto $\tau$ como $\tau'$ aumentan en $\Delta\tau$.

  \item Compare los resultados de los dos incisos. Explique por qué la respuesta del consumo depende de si el cambio en el ingreso disponible es transitorio o permanente. Relacione su respuesta con la hipótesis del ingreso permanente.
\end{enumerate}
\end{problema}



\begin{problema}{Efectos ingreso y sustitución de la tasa de interés real}

Para cada uno de los siguientes casos, utilice un diagrama en el plano $(c,c')$ para analizar los efectos de un aumento de la tasa de interés real. Grafique las restricciones presupuestarias inicial y final y construya una recta presupuestaria compensada para realizar una descomposición del cambio en la elección óptima en un efecto sustitución y un efecto ingreso. En cada caso, determine el signo del efecto sustitución, del efecto ingreso y del efecto total sobre $c$, $c'$ y $s$.

\begin{enumerate}
    \item Un hogar \textbf{ahorrador} para el cual, en la determinación de $c$, el efecto sustitución domina al efecto ingreso.

    \item Un hogar \textbf{ahorrador} para el cual, en la determinación de $c$, el efecto ingreso domina al efecto sustitución.

    \item Un hogar \textbf{deudor} para el cual, en la determinación de $c'$, el efecto sustitución domina al efecto ingreso.

    \item Un hogar \textbf{deudor} para el cual, en la determinación de $c'$, el efecto ingreso domina al efecto sustitución.
\end{enumerate}

\end{problema}






\begin{problema}{Hipótesis del ingreso permanente y función de consumo keynesiana}

Compare la función de consumo keynesiana con la hipótesis del ingreso permanente. En particular, explique:

\begin{enumerate}
    \item qué medida del ingreso determina el consumo presente en cada teoría;

    \item cómo responde el consumo corriente ante aumentos transitorios y permanentes del ingreso;

    \item cuál es el papel del ahorro y del endeudamiento en la determinación del consumo;

    \item qué evidencia empírica permitiría distinguir entre ambas explicaciones del comportamiento del consumo.
\end{enumerate}

\end{problema}






\begin{problema}{Estudio de caso}

Suponga que es un asesor del Secretario de Hacienda. El gobierno federal está diseñando un paquete fiscal para estimular la demanda agregada mediante un aumento de las transferencias y necesita proyectar sus efectos sobre el consumo y la producción.

Uno de los asesores estima la siguiente función de consumo:
\[
    C=1.04+0.54(Y-T+\mathrm{TR}),
\]
donde $Y-T+\mathrm{TR}$ es el ingreso disponible. Para realizar sus proyecciones, el asesor calcula el siguiente multiplicador fiscal:
    \[
        \mu_{\mathrm{TR}}
        =
        \frac{0.54}{1-0.54} \approx 1.17
    \]
Por lo tanto, proyecta que, por cada peso adicional de transferencias, la producción aumentará aproximadamente $1.17$ pesos.
\begin{enumerate}
    \item El asesor propone utilizar este multiplicador para proyectar los efectos de cualquier programa de transferencias. A la luz de la hipótesis del ingreso permanente, ¿considera apropiado utilizar estas estimaciones para realizar sus proyecciones? Argumente su respuesta.
\end{enumerate}

\end{problema}









\begin{problema}{Problema de dos periodos con riqueza inicial}
\label{Ej:P2P}
Suponga que un hogar resuelve el siguiente problema de consumo y ahorro en dos periodos:
\[
    \max_{c_1,s,c_2}\; u(c_1)+\beta u(c_2)
\]
sujeto a
\[
    s=s_0+y_1-\tau_1-c_1,
\]
\[
    c_2=y_2-\tau_2+(1+r)s,
\]

donde $c_1$ es el consumo en el periodo 1; $c_2$ es el consumo en el periodo 2; $y_1$ es el ingreso del hogar en el periodo 1; $y_2$ es su ingreso en el periodo 2; $\tau_1$ y $\tau_2$ son los impuestos pagados en los periodos 1 y 2, respectivamente; $s_0$ es la riqueza inicial del hogar; $s$ es el ahorro elegido en el primer periodo; $r$ es la tasa de interés real; y $\beta$ es el factor subjetivo de descuento.

Suponemos que el hogar tiene preferencias CRRA en cada periodo: 

\[
    u(c)=\frac{c^{1-\sigma}-1}{1-\sigma}, \qquad \sigma >0.
\]

\begin{enumerate}
    \item Obtenga la ecuación de Euler.
    \item Muestre que, bajo preferencias CRRA, la elasticidad de sustitución intertemporal es $1/\sigma$.\footnote{La elasticidad de sustitución intertemporal mide el cambio porcentual en la razón entre elección de consumo futuro y consumo presente ante un cambio porcentual en el rendimiento real. Es decir,
      \[
        \frac{\partial\log(c_2/c_1)}{\partial\log(1+r)}
      \].} Explique por qué un valor bajo de $1/\sigma$ implica una menor disposición a sustituir consumo entre periodos ante cambios en la tasa de interés real y, por tanto, una mayor preferencia por suavizar el consumo.

    \item Obtenga expresiones en forma cerrada para las elecciones óptimas del hogar, $c_1^*$, $c_2^*$ y $s^*$.\footnote{Una solución está en \textit{forma cerrada} cuando la variable de interés puede expresarse explícitamente como una función de otras variables y parámetros. Por ejemplo, si se solicita resolver para una variable $x$, una expresión como $x=y^2-z$ constituye una solución en forma cerrada. En cambio, una ecuación como $e^x+x-b=0$ determina implícitamente a $x$, pero no es una solución en forma cerrada.}

    \item ¿Cómo depende la razón $c_1^*/y_1$ del ingreso futuro $y_2$? ¿Qué ocurriría con el consumo presente si los hogares se volvieran repentinamente más optimistas acerca de su ingreso futuro? Interprete económicamente.

    \item ¿Cómo depende la razón $c_1^*/y_1$ de la riqueza inicial $s_0$? Interprete económicamente.

    \item ¿Cómo depende la razón $c_1^*/y_1$ del factor subjetivo de descuento $\beta$? Interprete económicamente.
\end{enumerate}

\end{problema}








\begin{problema}{Restricciones de crédito y la Equivalencia Ricardiana}

Suponga que un hogar resuelve la siguiente variante del problema del
Ejercicio~\ref{Ej:P2P}:
\[
    \max_{c_1,a,c_2}\;u(c_1)+\beta u(c_2)
\]
sujeto a
\[
    s=y_1-\tau_1-c_1,
\]
\[
    c_2=y_2-\tau_2+(1+r)s,
\]
y
\begin{equation}
    s\geq-b.
    \label{eq:restriccion-credito}
\end{equation}

\begin{enumerate}
    \item ¿Qué significa la restricción
    \eqref{eq:restriccion-credito}? ¿Qué representa $b$?

    \item Grafique la restricción presupuestaria del hogar. En la misma gráfica, represente la restricción \eqref{eq:restriccion-credito}.

    \item Resuelva las elecciones óptimas $c_1^*$, $c_2^*$ y $s^*$.

    \textit{Hint:} observe que la restricción
    \eqref{eq:restriccion-credito} es una desigualdad débil. Por lo tanto, la restricción puede saturarse o no saturarse. Si no se satura, puede utilizar la solución obtenida en el Ejercicio~\ref{Ej:P2P}. Después, determine qué ocurre cuando la restricción se satura y establezca las condiciones bajo las cuales se presenta cada caso.

    \item Muestre que, manteniendo todo lo demás constante, es más probable que la restricción \eqref{eq:restriccion-credito} se sature cuando:
    \begin{enumerate}
        \renewcommand{\labelenumii}{\textnormal{(\roman{enumii})}}
        \item el ingreso disponible futuro $y_2-\tau_2$ es alto;
        \item el ingreso disponible presente $y_1-\tau_1$ es bajo;
        \item el límite de endeudamiento $b$ es bajo.
    \end{enumerate}
    Interprete cada una de estas condiciones.

    \item Suponga que el gobierno anuncia un paquete de estímulo fiscal de tamaño $\Delta$. El programa consiste en reducir $\tau_1$ en $\Delta$ y aumentar $\tau_2$ en $\Delta(1+r)$, de modo que el valor presente de los impuestos permanece constante.

    ¿Cómo responde $c_1$ al paquete de estímulo si la economía parte de una situación en la que la restricción \eqref{eq:restriccion-credito} no se satura? ¿Cómo responde $c_1$ si la economía parte de una situación en la que la restricción sí se satura? Explique.
\end{enumerate}

\end{problema}





\end{document}
